% ============================================================
% Review Paper: Cost–Benefit Analysis and Literature Review
% of Industrial Load Flexibility
% ============================================================
% Overleaf-ready template — upload this .tex file directly
% Compile with pdflatex or lualatex
% ============================================================

\documentclass[12pt, a4paper, twocolumn]{article}

% ---- Packages ----
\usepackage[utf8]{inputenc}
\usepackage[T1]{fontenc}
\usepackage{lmodern}
\usepackage[english]{babel}
\usepackage{geometry}
\geometry{margin=1in}
\usepackage{graphicx}
\usepackage{booktabs}
\usepackage{longtable}
\usepackage{array}
\usepackage{multirow}
\usepackage{tabularx}
\usepackage{amsmath,amssymb}
\usepackage{hyperref}
\hypersetup{
    colorlinks=true,
    linkcolor=blue,
    citecolor=blue,
    urlcolor=blue
}
\usepackage{natbib}
\bibliographystyle{plainnat}
\usepackage{enumitem}
\usepackage{float}
\usepackage{caption}
\usepackage{subcaption}
\usepackage{xcolor}
\usepackage{soul}              % for highlighting / markup
\usepackage{lineno}            % line numbers for review
\usepackage{setspace}
\usepackage{titlesec}

% ---- Formatting ----
\titleformat{\section}{\large\bfseries}{\thesection.}{0.5em}{}
\titleformat{\subsection}{\normalsize\bfseries}{\thesubsection}{0.5em}{}
\titleformat{\subsubsection}{\normalsize\itshape}{\thesubsubsection}{0.5em}{}

% Uncomment below for line numbers during review
% \linenumbers

% ---- Custom commands ----
\newcommand{\LF}{\textit{LF}}
\newcommand{\kW}{\,\text{kW}}
\newcommand{\kWh}{\,\text{kWh}}
\newcommand{\USD}{\$}
\newcommand{\TODO}[1]{\textcolor{red}{\textbf{[TODO: #1]}}}

% ============================================================
\begin{document}

% ---- Title Block ----
\title{%
  \textbf{Cost--Benefit Analysis and Literature Review of\\
  Industrial Load Flexibility:\\
  Financial Value, Stakeholder Impact, and\\
  Utility Rebate Program Inventory}}

\author{
  Author One\thanks{Affiliation, e-mail}
  \and
  Author Two\thanks{Affiliation, e-mail}
  \and
  Author Three\thanks{Affiliation, e-mail}
}

\date{\today}

\twocolumn[
  \begin{@twocolumnfalse}
    \maketitle
    % ---- Abstract ----
    \begin{abstract}
    \noindent
    Industrial load flexibility (\LF{}) — the ability of a facility's electrical demand to be shifted, shed, or modulated in response to economic, operational, or grid signals — is increasingly recognised as a valuable asset for large commercial and industrial (C\&I) consumers.  Despite its significance, no comprehensive review has simultaneously addressed (i)~the direct financial benefits of \LF{}, (ii)~the indirect and non-financial value created for internal stakeholders who may not directly benefit from utility bill reduction, and (iii)~the landscape of utility incentive and rebate programmes that reward \LF{}.  This paper fills that gap by synthesising cost--benefit evidence from the academic literature, industry case studies, and regulatory filings.  We further present an inventory of \LF{}-linked rebate programmes across major U.S.\ and international utility territories, documenting incentive magnitudes, calculation methodologies, and eligibility criteria.  Our findings indicate that \TODO{summarise key findings here}.  We conclude with a research agenda for quantifying the full value chain of \LF{} in an era of increasing electrification and behind-the-meter distributed energy resources.

    \vspace{0.5em}
    \noindent\textbf{Keywords:} load flexibility, demand management, industrial energy efficiency, utility rebate programmes, cost--benefit analysis, demand response, peak demand reduction, load shifting
    \end{abstract}
    \vspace{1em}
  \end{@twocolumnfalse}
]

% ============================================================
% TABLE OF CONTENTS (optional — remove for final submission)
% ============================================================
% \tableofcontents
% \newpage

% ============================================================
\section{Introduction}
\label{sec:introduction}
% ============================================================

\subsection{Background and Motivation}
\TODO{Introduce the concept of load flexibility and why it matters for industrial electricity consumers.}

\begin{itemize}[nosep]
  \item Definition of load flexibility: the capacity of an industrial facility to adjust the timing, magnitude, or pattern of its electricity consumption in response to price signals, grid conditions, or operational objectives.
  \item Historical context: evolution of demand-based tariff structures and demand-response programmes.
  \item Growing relevance due to electrification of industrial processes, EV fleet charging, on-site generation, battery storage, and grid decarbonisation.
\end{itemize}

\subsection{Problem Statement}
\TODO{Articulate the research gap.}

\begin{itemize}[nosep]
  \item Lack of a unified review combining direct financial, indirect financial, and non-financial value of \LF{}.
  \item No comprehensive catalogue of utility \LF{} rebate/incentive programmes.
  \item Need for actionable data to support capital-investment decisions and internal business cases.
\end{itemize}

\subsection{Research Questions}
\begin{enumerate}[nosep]
  \item What is the documented direct financial value (cost savings) of \LF{} for industrial consumers?
  \item What indirect and non-financial benefits does \LF{} confer on internal stakeholders (e.g., operations, maintenance, sustainability, finance) who may not directly benefit from utility bill reduction?
  \item What \LF{}-linked rebate and incentive programmes exist across utility territories, and how do they vary in magnitude, eligibility, and calculation methodology?
\end{enumerate}

\subsection{Scope and Delimitations}
\TODO{Define geographic scope, industry sectors covered, date range of literature, exclusion criteria.}

\subsection{Paper Organisation}
\TODO{Briefly outline the structure of the remaining sections.}


% ============================================================
\section{Methodology}
\label{sec:methodology}
% ============================================================

\subsection{Literature Search Strategy}
\begin{itemize}[nosep]
  \item Databases searched (e.g., Scopus, Web of Science, IEEE Xplore, Google Scholar, DOE EERE, ACEEE).
  \item Search terms and Boolean queries.
  \item Date range and language filters.
  \item Grey literature: utility tariff books, regulatory dockets, trade publications, white papers.
\end{itemize}

\subsection{Inclusion and Exclusion Criteria}
\TODO{Define criteria for including/excluding studies.}

\begin{table}[H]
\centering
\caption{Inclusion / Exclusion Criteria}
\label{tab:criteria}
\begin{tabularx}{\columnwidth}{lX}
\toprule
\textbf{Criterion} & \textbf{Description} \\
\midrule
Population   & Industrial / large C\&I electricity consumers \\
Intervention & \LF{} measures, enabling technologies, or tariff designs \\
Outcome      & Financial savings, non-financial benefits, or rebate data \\
Language     & English \\
Date range   & \TODO{e.g., 2000--2026} \\
\bottomrule
\end{tabularx}
\end{table}

\subsection{Data Extraction and Synthesis}
\TODO{Describe how data were extracted, coded, and synthesised (narrative synthesis, meta-analysis, etc.).}

\subsection{Utility Rebate Programme Inventory Method}
\begin{itemize}[nosep]
  \item Sources: utility tariff schedules, state PUC filings, DSIRE database, direct utility enquiry.
  \item Data fields captured: utility name, territory, rate class, \LF{} threshold, incentive amount (\$/kW or \$/kWh), calculation formula, eligibility requirements, programme term.
\end{itemize}


% ============================================================
\section{Fundamentals of Industrial Load Flexibility}
\label{sec:fundamentals}
% ============================================================

\subsection{Definition and Characterisation}
Load flexibility (\LF{}) can be characterised along several dimensions:
\begin{itemize}[nosep]
  \item \textbf{Magnitude} — the amount of demand (kW) that can be shifted, shed, or modulated.
  \item \textbf{Duration} — how long a load adjustment can be sustained.
  \item \textbf{Response time} — latency between a signal and realised load change.
  \item \textbf{Recovery} — any rebound or payback load following a flexibility event.
  \item \textbf{Frequency} — how often flexibility can be called upon within a period.
\end{itemize}

A common summary metric relates average demand to peak demand:
\begin{equation}
  \text{Load Factor} = \frac{\displaystyle\sum_{i=1}^{N} P_i \cdot \Delta t}{P_{\max} \cdot T} \times 100\%
  \label{eq:lf}
\end{equation}
where higher values indicate a flatter, more flexible load profile.  However, \LF{} encompasses capabilities beyond this single ratio.

\subsection{Barriers to Load Flexibility in Industrial Facilities}
\begin{itemize}[nosep]
  \item Rigid batch-process scheduling and shift patterns.
  \item Seasonal or weather-driven load variability.
  \item Start-up transients and coincident equipment operation.
  \item Lack of visibility into real-time demand and limited sub-metering.
  \item Organisational silos between energy management and production.
\end{itemize}

\subsection{Enabling Strategies and Technologies}
\begin{itemize}[nosep]
  \item Load shifting and scheduling optimisation.
  \item Peak shaving via battery energy storage systems (BESS).
  \item On-site generation dispatch and co-generation.
  \item Demand-controlled ventilation and process staging.
  \item Advanced metering, SCADA integration, and automated demand-response platforms.
  \item Thermal energy storage for HVAC and process loads.
\end{itemize}


% ============================================================
\section{Direct Financial Value of Load Flexibility}
\label{sec:direct-value}
% ============================================================

\subsection{Impact on Demand Charges}
\TODO{Review studies quantifying demand-charge savings from \LF{}.  Include representative ranges (\$/kW-month) by rate class.}

\subsection{Impact on Energy Rates and Ratchet Clauses}
\TODO{Discuss how limited \LF{} interacts with ratchet demand clauses, time-of-use rates, and real-time pricing.}

\subsection{Impact on Transmission and Distribution Costs}
\TODO{Review evidence on how \LF{} affects allocated T\&D costs under cost-of-service regulation.}

\subsection{Total Cost-of-Electricity Sensitivity Analysis}
\TODO{Present or summarise published models showing total \$/kWh as a function of \LF{} capability for representative industrial tariffs.}

\begin{table}[H]
\centering
\caption{Representative Demand Charge Savings from \LF{}}
\label{tab:savings}
\begin{tabularx}{\columnwidth}{lccX}
\toprule
\textbf{Source} & \textbf{Baseline \LF{}} & \textbf{Improved \LF{}} & \textbf{Annual Savings} \\
\midrule
\TODO{} & \TODO{} & \TODO{} & \TODO{} \\
\TODO{} & \TODO{} & \TODO{} & \TODO{} \\
\TODO{} & \TODO{} & \TODO{} & \TODO{} \\
\bottomrule
\end{tabularx}
\end{table}

\subsection{Case Studies}
\TODO{Summarise 3--5 published industrial case studies with quantified savings.}


% ============================================================
\section{Indirect and Non-Financial Value of Load Flexibility}
\label{sec:indirect-value}
% ============================================================

\subsection{Operational Reliability and Equipment Longevity}
\TODO{Discuss how smoother load profiles reduce thermal cycling, mechanical stress, and unplanned downtime.}

\subsection{Power Quality and Process Stability}
\TODO{Link \LF{} to voltage stability, reduced flicker, and process yield.}

\subsection{Grid Services and Demand Response Eligibility}
\TODO{Explain how facilities with high \LF{} can participate in DR and ancillary-service markets, generating additional revenue streams.}

\subsection{Sustainability and Carbon Reduction}
\TODO{Review evidence that peak demand correlates with marginal emissions; quantify carbon value.}

\subsection{Internal Stakeholder Value Mapping}
\begin{table}[H]
\centering
\caption{Non-Financial Value of \LF{} by Internal Stakeholder Group}
\label{tab:stakeholder}
\begin{tabularx}{\columnwidth}{lX}
\toprule
\textbf{Stakeholder} & \textbf{Value Created} \\
\midrule
Operations     & Reduced peak-related curtailments; smoother scheduling; production continuity \\
Maintenance    & Lower thermal stress; extended transformer/motor life \\
Finance / Treasury & Predictable utility spend; reduced budget variance \\
Sustainability / ESG & Lower Scope 2 emissions intensity; reporting metrics \\
Procurement    & Stronger position in retail energy negotiations \\
Capital Projects & Deferred substation/switchgear upgrades \\
Risk Management & Reduced exposure to demand-ratchet penalties \\
\bottomrule
\end{tabularx}
\end{table}

\subsection{Organisational and Behavioural Factors}
\TODO{Review literature on awareness, training, and incentive alignment that influences \LF{} outcomes.}


% ============================================================
\section{Inventory of Utility Load Flexibility Rebate Programmes}
\label{sec:rebate-inventory}
% ============================================================

\subsection{Overview of Programme Landscape}
\TODO{Summarise how many utilities offer \LF{}-linked incentives, geographic distribution, and trends over time.}

\subsection{Programme Typology}
\begin{enumerate}[nosep]
  \item \textbf{Tariff-embedded \LF{} adjustments} — rate discount/premium tied to demonstrated \LF{}.
  \item \textbf{Demand-side management (DSM) rebates} — capital incentives for \LF{}-enabling equipment.
  \item \textbf{Custom/calculated incentives} — utility-specific formulas based on verified kW reduction or load-shape improvement.
  \item \textbf{Performance-based incentives} — ongoing payments for sustained \LF{} above a threshold or participation in flexibility events.
\end{enumerate}

\subsection{Incentive Calculation Methods}
\TODO{Describe common formulas (e.g., \$/kW-saved, percentage discount on demand charge, sliding-scale adjustments).}

\subsection{Eligibility Requirements}
\TODO{Summarise typical requirements: rate class, minimum demand, metering, M\&V protocols, pre-approval.}

\subsection{Detailed Programme Inventory}

\begin{table*}[ht]
\centering
\caption{Inventory of Utility \LF{} Rebate and Incentive Programmes (Representative Sample)}
\label{tab:inventory}
\small
\begin{tabularx}{\textwidth}{llllXll}
\toprule
\textbf{Utility / Territory} & \textbf{State} & \textbf{Rate Class} & \textbf{LF Threshold} & \textbf{Incentive Description} & \textbf{Amount} & \textbf{Method} \\
\midrule
\TODO{Utility A} & \TODO{} & \TODO{} & \TODO{} & \TODO{} & \TODO{} & \TODO{} \\
\TODO{Utility B} & \TODO{} & \TODO{} & \TODO{} & \TODO{} & \TODO{} & \TODO{} \\
\TODO{Utility C} & \TODO{} & \TODO{} & \TODO{} & \TODO{} & \TODO{} & \TODO{} \\
\TODO{Utility D} & \TODO{} & \TODO{} & \TODO{} & \TODO{} & \TODO{} & \TODO{} \\
\TODO{Utility E} & \TODO{} & \TODO{} & \TODO{} & \TODO{} & \TODO{} & \TODO{} \\
\bottomrule
\end{tabularx}
\end{table*}

\TODO{Full inventory to be provided in supplementary material / appendix.}

\subsection{Comparative Analysis of Programmes}
\TODO{Analyse variation in incentive generosity, threshold stringency, and programme maturity across territories.}


% ============================================================
\section{Synthesis and Discussion}
\label{sec:discussion}
% ============================================================

\subsection{Key Findings}
\TODO{Synthesise main findings across the three review dimensions.}

\subsection{Gaps in the Literature}
\begin{itemize}[nosep]
  \item Limited peer-reviewed studies quantifying non-financial \LF{} value.
  \item Lack of standardised metrics for comparing rebate programme effectiveness.
  \item Sparse data on \LF{} interactions with behind-the-meter DERs and storage.
  \item \TODO{Additional gaps identified during review.}
\end{itemize}

\subsection{Implications for Industrial Energy Managers}
\TODO{Translate findings into actionable guidance.}

\subsection{Implications for Utility Programme Design}
\TODO{Recommendations for utilities designing or revising \LF{} incentive structures.}

\subsection{Limitations of This Review}
\TODO{Discuss limitations: publication bias, geographic skew, data currency, etc.}


% ============================================================
\section{Future Research Directions}
\label{sec:future}
% ============================================================

\begin{enumerate}[nosep]
  \item Development of a standardised \LF{} value framework integrating direct, indirect, and non-financial benefits.
  \item Empirical studies on the interaction of \LF{} with on-site BESS, solar PV, and other DERs.
  \item Longitudinal analysis of rebate programme participation and verified savings.
  \item Machine-learning approaches to real-time \LF{} forecasting and automated demand control.
  \item Cross-jurisdictional policy analysis of \LF{} incentive design best practices.
\end{enumerate}


% ============================================================
\section{Conclusion}
\label{sec:conclusion}
% ============================================================

\TODO{Summarise the review's contributions, restate key findings, and close with a call to action for researchers and practitioners.}


% ============================================================
% REFERENCES
% ============================================================
% Use BibTeX: create a file called references.bib and
% uncomment the line below.  Placeholder entries are included
% inline for convenience during drafting.
% \bibliography{references}

\begin{thebibliography}{99}

\bibitem{doe2020}
U.S.\ Department of Energy, Office of Energy Efficiency \& Renewable Energy.
\textit{Industrial Assessment Centers Database}.
\url{https://iac.university}, 2020.

\bibitem{eia2023}
U.S.\ Energy Information Administration.
\textit{Annual Electric Power Industry Report (EIA-861)}.
2023.

\bibitem{aceee2019}
American Council for an Energy-Efficient Economy.
\textit{The State of Industrial Energy Efficiency Programs}.
Research Report, 2019.

\bibitem{dsire2024}
Database of State Incentives for Renewables \& Efficiency (DSIRE).
\url{https://www.dsireusa.org/}, Accessed 2024.

\TODO{Add remaining references.}

\end{thebibliography}


% ============================================================
% APPENDICES
% ============================================================
\appendix

\section{Full Utility \LF{} Rebate Programme Inventory}
\label{app:inventory}
\TODO{Include the complete programme-by-programme inventory table here or as supplementary material.}

\section{Search Strategy Details}
\label{app:search}
\TODO{Full Boolean search strings, database-specific syntax, and PRISMA flow diagram.}

\section{Data Extraction Template}
\label{app:extraction}
\TODO{Blank data extraction form used for each included study.}

% ============================================================
\end{document}
